\documentclass[10pt]{article}
\usepackage[margin=20mm]{geometry}
\usepackage{amsmath,amssymb,fontspec,unicode-math,xurl,hyperref}
\IfFontExistsTF{DejaVu Serif}{\setmainfont{DejaVu Serif}}{\setmainfont{Latin Modern Roman}}
\setmathfont{latinmodern-math.otf}
\setlength{\parindent}{0pt}\setlength{\parskip}{6pt}\sloppy
\begin{document}
0076 | Scientific verification report\par
CAT Merging: A Training-Free Approach for Resolving Conflicts in Model Merging\par
I25-43445 | ICML 2025 | Review date: 24 October 2026\par
Batch B001 / [internal task] | [internal identifier omitted]\par
\# 1. Identity and reviewed version\par
Authors: Wenju Sun, Qingyong Li, Yangli-ao Geng, and Boyang Li. Reviewed source: the publisher-linked PMLR volume 267 PDF, 21 pages including Appendices A-C. Physical PDF pages 1-21 match the printed article numbers; proceedings pages 57523-57543 are not used as locators.\par
The preserved file is `papers/0076/0076.pdf`. SHA-256 a27859400d498e67f56791d8276bd024d00f8ef78ba89ae13b7dc50c274fa432 matches `verification/0076/controller\_acquisition\_receipt.json`. Title, authors, and page count were checked. Source bytes were not replaced.\par
The public author repository was inspected at commit ab87731e405b9eeca1f822685ede0276d756907c, dated 16 June 2026. Its CAT entry point and relevant shared source files are preserved under `sources/0076/model-merging`. This is a later released snapshot, not an authenticated record of the 2025 experimental runs.\par
\# 2. Outcome and checked scope\par
Supported impact: Minor. Static inspection and an exact independent calculation demonstrate that the checked release selects the wrong spectral directions in its fixed-rank branch. The default entry point also disables evaluation. These findings affect the supplied implementation; the record does not establish which code generated the paper’s tables. Separately, three printed row means are inconsistent with their printed inputs but do not change the central CAT rankings.\par
No substantive mathematical defect was found in Lemma 4.2, Theorem 4.4, or the constrained local trimming derivations under their stated and contextually explicit assumptions. Main Sections 3-6, Appendix A, numerical tables, experimental descriptions, and the active released implementation were checked. Independent calculations did not execute or import author code. Benchmark accuracy, runtime, sample selection, and training provenance were not reproduced.\par
This report separates the validity of the paper’s local mathematics, a demonstrated later-release implementation discrepancy, and the unresolved provenance of the published experiments. Scientific decisions and human verification remain pending.\par
\# 3. Definitions, assumptions, and reasoning that survived\par
The task vector is T\_i=W\_i-W\_0. At a linear layer, features X\_k are held fixed while the paper trims another task vector by Phi\_k(T\_i)=T\_i-T\_i B\_k B\_k\textasciicircum{}T. The local objective adds cross-task feature disturbance to a lambda-weighted own-task feature deviation. The basis has c orthonormal columns; c is fixed by the method.\par
On pp.3 and 13-14, the triangle inequality splits propagated and local feature shifts. The stated layer Lipschitz constant must apply to the compared input feature pair at the perturbed layer weights. With the original input unchanged, the initial shift is zero; recursion and the final loss Lipschitz bound give Theorem 4.4. This is a local perturbation bound, not a theorem of globally optimal sequential merging or monotone test accuracy.\par
On pp.4 and 14-15, P=B B\textasciicircum{}T is an orthogonal projector. Orthogonal components have zero Frobenius inner product, so the norm decomposition is exact. The objective becomes a constant minus tr(B\textasciicircum{}T G B), where G=sum\_\{i != k\} T\_i\textasciicircum{}T(X\_k\textasciicircum{}T X\_k-lambda X\_i\textasciicircum{}T X\_i)T\_i. The largest c algebraic eigenvalues maximize the trace even if G is indefinite.\par
On pp.5 and 15-16, the mask selection is correct when read with the explicit top-c cardinality stated in the prose. Appendix A.6 additionally assumes equal sample counts and 0<lambda<1, making n(1-lambda) a positive common factor. Its conclusion must not be extended to lambda>1 without revisiting the shift objective. Figure 4 explores larger lambda values; that sensitivity experiment is not a proof of the narrower shift lemma outside its assumptions.\par
A stray x-dependent intermediate row in Eq.(30) can be removed while preserving the direct valid equality and conclusion. It is a presentation issue outside this task’s substantive-error scope. Noncommutation of successive orthogonal projectors alone is also not evidence of an increase in the fixed-input Frobenius disturbance.\par
\[J(B)=\sum_{i\ne k}\left(\|X_kT_i(I-BB^\top)\|_F^2+\lambda\|X_iT_iBB^\top\|_F^2\right)=C-\operatorname{tr}(B^\top GB),\quad B^\top B=I_c.\]\par
\# 4. Problem 1: the checked release does not implement the paper’s spectral optimizer\par
Source claim and location. Section 5.1, Eq.(10)-(11), p.4, and Appendix A.4, pp.14-15, require the top-c eigenvectors of G. The released methods/cat\_merging\_utils.py constructs G with the same sign, calls torch.linalg.eigh at line 176, and selects eigenvectors[:,0:r] at line 189. PyTorch documents ascending eigenvalue order. Thus the selected fixed-rank subspace contains the smallest, not the largest, eigenvalues. The active entry point sets ratio=4.\par
Complete independent check. Set K=2, d=h=n=2, T\_i=I\_2, X\_i=I\_2, lambda=1/2, X\_k=diag(sqrt(3/2),sqrt(9/2)), and c=1. These are real, dimensionally compatible feature/task matrices with equal sample counts and a positive permitted lambda. Both e\_1 and e\_2 are feasible unit-norm bases. G=diag(1,4), so the positive-eigenvalue branch is active. The constant C is 6. The code’s e\_1 gives J=6-1=5; the paper’s e\_2 gives J=6-4=2. Direct calculation gives the same values: 9/2+1/2 versus 3/2+1/2. An exact rational five-dimensional c=4 analogue also gives 15/2 versus 7/2, so the discrepancy is not confined to c=1.\par
Alternative readings and counterevidence. This is not a transpose or sign convention: the inspected G matches the paper after accounting for PyTorch’s stored weight orientation. The ratio=-1 branch selects all positive eigenvectors and is not alleged to have this ordering defect. The comparison isolates basis selection before the release’s later norm rescaling at lines 224-228; that rescaling is itself a different operation from the pure projection in the paper and does not validate the claimed fixed-rank optimizer.\par
Impact and severity. The released core trimming step can fail to minimize the paper’s exact local objective. The paper’s derivation remains correct. Because the inspected commit postdates publication and no run-to-commit mapping was obtained, this is a confirmed Minor release discrepancy, with impact on published experiments marked Needs Evidence. A Major finding against the published scientific results would require the missing dependency evidence, not merely the presence of this later code.\par
Scoped repair. For the fixed-c branch, 1<=c<=d, selecting eigenvectors[:,-c:] restores the top-c spectral choice. This is analytically verified for that step and confirmed by the independent example. It does not validate the later rescaling, the all-negative-spectrum guard, the complete pipeline, or benchmark results. No author source was changed and no adoption is claimed.\par
Additional release limitation. main\_catmerging.py lines 78-80 set ratio\_ln=ratio\_bias=0 and test\_datasets=[]. Consequently lines 109-114 perform no dataset evaluation, and lines 127-128 subsequently index empty summaries if execution reaches them. An earlier or differently configured experiment script could have been used. Restoring those settings requires an authenticated experimental configuration; the report does not invent one.\par
\# 5. Problem 2: three printed means do not agree with the printed rows\par
The checked quantities are ordinary unweighted row means in Tables 1-3, pp.6-7. Independent standard-library arithmetic recomputed 58 displayed row averages, all ten Table 10 balance statistics, and the best-count columns. Three discrepancies exceed the uncertainty from rounding every input cell to its displayed precision.\par
Table 1 Traditional MTL: printed 88.9 versus arithmetic 89.0625 from the eight displayed values. Table 2 Traditional MTL: printed 93.5 versus 93.0875. Table 3 TATR: printed 58.39 versus 58.7525. For one-decimal inputs, any mean of the underlying rounded observations can differ from the displayed-input mean by at most 0.05; for two-decimal inputs the bound is 0.005. The displayed means lie outside those ranges even allowing their own final rounding.\par
Counterevidence and impact. All three discrepancies concern comparator summaries; none changes the reported CAT ranking or the main vision gains. The central CAT/PCB means and differences support 2.5 and 2.0 percentage points. Table 10 matches sample standard deviations (denominator n-1), including CAT values 6.2112656865 and 2.5116514203, and supports the stated balance ordering. The aggregate discrepancy is Minor and does not support a claim of benchmark failure or intentional misreporting.\par
A precise conditional correction is available: if the printed component scores are authoritative and equal task weighting was intended, the three means round to 89.1, 93.1, and 58.75. This repairs displayed aggregation only. Unrounded logs would determine whether a component entry instead needs correction; those logs were not available.\par
\# 6. Dependence, versions, and remaining evidence\par
The theory supports a constrained local surrogate; empirical superiority remains an experimental claim. The code discrepancy changes an implementation of that surrogate, but its connection to the 2025 experiments is unresolved. Required follow-up is the experimental commit, configuration, checkpoints and run logs, followed by appropriately scoped fresh experiments if those records implicate the defective branch.\par
The retrieved arXiv record identifies v2 dated 14 May 2025 and an abstract gain of 2.5 percentage points, consistent with the preserved PDF. The current publisher landing page instead lists 4.7\% and omits language tasks in the abstract. That metadata/version discrepancy is recorded; the review uses the preserved PDF and does not silently substitute the landing-page abstract.\par
No later correction to the inspected claims was located in the bounded publisher/arXiv/repository checks. This is not an exhaustive no-errata certificate. The OpenReview web page returned a browser-verification challenge. NLP metric provenance, checkpoint-level identity, uncertainty over independent experimental runs, the timing measurements, and the plotted loss surfaces remain unverified.\par
\# 7. Conclusion and repair status\par
Item | Judgment / scope | Repair verification | Adoption / remaining work\par
Spectral selection and release entry point | Minor, checked 2026 release; published-run link Needs Evidence | Largest-c slice analytically valid for fixed-c step only; no complete pipeline repair | Not requested; obtain original experiment commit/configuration\par
Three printed means | Minor local numerical aggregation discrepancy | Conditional mean corrections independently computed | Not requested; compare unrounded result logs\par
Lemma 4.2, Theorem 4.4, local trimming theory | No Issue Found within checked assumptions; scope clarification retained | No substantive repair proposed | Human review pending\par
Accuracy, runtime and plotted surfaces | Not empirically reproduced | No reproduction claim | Fresh experiments only if separately authorized\par
\# 8. Consulted sources and stable identifiers\par
R1. Sun et al. (2025), CAT Merging. PMLR 267:57523-57543. Primary preserved PDF and included appendices; SHA-256 stated in Section 1. Publisher: https://proceedings.mlr.press/v267/sun25i.html\par
R2. SunWenJu123/model-merging, commit ab87731e405b9eeca1f822685ede0276d756907c (16 June 2026), CAT entry point, active merging implementation, evaluation utilities and dataset support. https://github.com/SunWenJu123/model-merging/tree/ab87731e405b9eeca1f822685ede0276d756907c\par
R3. PyTorch 2.14 documentation, torch.linalg.eigh, ascending eigenvalue order and corresponding eigenvector columns; accessed 24 October 2026. https://docs.pytorch.org/docs/2.14/generated/torch.linalg.eigh.html\par
R4. arXiv:2505.06977v2, version metadata (14 May 2025); metadata cross-check, not a replacement for the primary PDF. https://arxiv.org/abs/2505.06977\par
\# Appendix A. Evidence, replay, and provenance\par
The machine-readable coverage, findings and provenance records are `verification/0076/coverage.json`, findings.json and provenance.json. The acquisition receipt, reviewer\_identity\_check.json, extraction\_manifest.json, math\_review\_notes.json, code\_review\_notes.json, empirical notes, code\_line\_extracts.txt and source snapshots preserve the checked evidence.\par
Exact finite replay: `verification/0076/code\_counterexample\_check.py` and code\_counterexample\_result.json. Table replay: `verification/0076/empirical\_replay.py` and empirical\_results.json. These are independently written reviewer calculations and do not import or execute author code. Their outputs, assumptions and actual results are preserved.\par
Text was extracted using the existing pdftotext executable with explicit physical page ranges. Seven mathematical pages (3-5 and 13-16) were rendered with existing pdftoppm and visually checked; the primary source page 4 was also independently inspected by the coordinating AI reviewer. Full-text control glyphs made a naive form-feed split produce 45 false pages; that failed split was retained and replaced with 21 explicit per-page extractions.\par
Tool limitations/failures: Lumen’s direct PDF reader returned Cannot read properties of undefined (reading ’toLowerCase’); the existing Poppler tools supplied the successful alternative. No software was installed. This replacement PDF was compiled offline from the accompanying TeX during acceptance; the controller remains responsible for promotion and final file acceptance.\par
Roles: the coordinating ChatGPT reviewer and three AI subagents checked mathematical reasoning, static code, and numerical evidence. Their agreement is not human verification. The user/controller retains all scientific decisions. Human review: pending. Author contact/adoption: not requested. No publication, model/API spending, author-code execution, queue edit, or controller-file mutation occurred.\par
\end{document}
